\documentclass[final,a4paper,12pt]{article}

\usepackage{notes}
\setlength\parskip{1ex}
\setlength\parindent{0cm}
\usepackage{float}
\usepackage{mathtools}

\newcommand\qeq{\stackrel{\mathclap{\normalfont\mbox{?}}}{=}}

\title{Geophysics III}
\author{Derivations Tutorial}
\date{}

\begin{document}
\maketitle

\section{Introduction}
{\it What are derivations and why do we do them?} \\
A derivation is the process of developing a relationship between one or more
variables that provide physical insight into a scientific problem.  These can be
quite formal, but are often used to develop a mathematical formulation to
predict or model observations.  Derivations can also be used to develop
corrections to data (you'll see the end result in the gravity corrections).

{\it Why do I make you do derivations?} \\
Some of you will enjoy them, but for the rest, I do have very good reasons for
having you perform derivations.  I can give you a fish, or teach you how to
catch one.  If you learn how to do a derivation it will give you the experience
that will help you follow along when you read papers involving derivations and
perhaps help you gain confidence to do them in the future if the need arise.

The other reason I have you do derivations is that it will help you gain
insights into a problem that you otherwise might not get from a single equation.
Often a derivation starts from a relatively general point and through a series
of approximations and/or constraints arrives at a specific solution.  The
solution may not be true in all cases as a result of these assumptions and
constraints.  Seeing, rather than being told, why the solution is limited is an
important part of understanding a physical process.  I want you to gain the
fundamental insight that may be useful in the future, while the equations we
derive may not.

\section{How to do a derivation}
For this course, I want you to be as organized and prescriptive as possible when
setting up and then deriving your solutions.  Being methodical will help make
the process of easier.  It also makes derivations easier to read---and easier to
assign partial credit!  It may not be necessary to do each of these steps every
time, but these are the generalized steps and reasons why they are a good idea.
There are 4 steps.  To remember them you can use the acronym: {\bf D$^{3}$C},
which stands for {\bf d}raw, {\bf d}efine, {\bf d}erive, and {\bf c}heck.  {\bf
Do NOT input values before you complete the derivation.}
\begin{enumerate}
    \item {\bf Draw} a picture and annotate the drawing with the symbols you
    will use in the derivation.

    Drawing a picture can help focus the problem, but also gives you an
    opportunity to see all the parts and how they relate to one another.

    \item {\bf Define} the symbols and give their values and units, including
    physical constants.  You can do this as a table or as a simple list.  This
    is also a good spot to identify the key assumptions and/or constraints that
    you may know {\it a priori}.  See how I have done the example below.  You
    might leave a few lines space below the list in the event that you need to
    add variables at a later point.

    Making a list of all the symbols and their meaning improves the readability
    in case you use a different symbol than someone else might for the same
    variable.  It also gives you the opportunity to identify the knowns and
    unknowns which can help as you decide what you need to do to solve the
    problem.
    
    It is best to use symbols by convention of the field within which
    they are used so that it avoids confusion.

    \item {\bf Derive} a key result, taking care to keep the equations in
    symbolic notation (don't plug in values).  Enumerate key equations that you
    know you will need up front.  This includes any systems of equations,
    initial and boundary conditions.  Number these equations as it will make it
    easier to refer to them later.  Identify the key steps, locations where:
    assumptions/constraints are applied; approximations are made; and a shift in
    focus is made (e.g. you start deriving a secondary result before
    continuing).  While you are working up your derivation, think about where
    you where you want to go and what you need to do to get there.

    This not only makes it easier to assign partial credit, but also means that
    the solution you have at the end can be used for any set of inputs, not just
    the ones you currently have.  If we always inserted values at the earliest
    possible point, we would have to rederive the result every time we wanted to
    change our inputs.

    \item {\bf Check} your answer.  Once you have a symbolic solution YOU ARE
    NOT DONE.  Since you may not know what you are looking for, there are
    several ways to check your solution.  First, perform a dimensional analysis
    (check units) if the units work out properly you've cleared the simplest
    hurdle.  Next, make sure the solution satisfies the initial and boundary
    conditions.  If you know the solution is asymptotic, make sure it satisfies
    those limits (e.g., if the result should approach zero as position or time
    approaches infinity).  Finally, you can graph the result as a final test, but
    generally the first two are sufficient.

    Does the solution make sense.  Generally, if the solution satisfies the
    applied constraints and the units work out properly you likely have the
    correct solution to the problem.  If you do the test and it doesn't satisfy
    the constraints or dimensional analysis, make a note of it and try to
    identify the location where things most likely went awry.  Doing so will
    vastly improve your partial credit when things don't go as planned.

    Checking your answer is another way to get partial credit if your solution
    isn't correct.  If you can't figure out the answer to a problem, but check
    the solution and can articulate why your answer cannot possibly be correct
    you will earn some points you would have otherwise lost.

    \item Once you are satisfied with your result, input values if necessary and
    circle your solution (it helps me know when you think you are done).  I also
    encourage you to discuss your derivations with someone else since it is one
    of the best ways to make sure that it is clear and complete.  Try them
    yourself, but I encourage you to work in groups to find the solutions.
\end{enumerate}


\section{Example}
We will be doing similar isostasy (buoyancy related) problems in a couple weeks
so don't worry if you don't understand the entire derivation.  You will soon.

\noindent {\bf Problem:} \\
What is the height mantle root necessary for an Airy model of isostatic
compensation?  Write in terms of the density contrast between the crust and
mantle (i.e.\ $\Delta \rho = \rho_{m} - \rho_{c}$).
\clearpage

\noindent {\bf Answer:} \\
{\bf Draw}\\
We begin by drawing a picture of the system and labeling the important
components.
\begin{figure}[H]
  \centerline{
    \includegraphics[]{columns.eps}
  }
  \caption{Isostatic columns.}
  \label{fig-columns}
\end{figure}
Note that the figure is cartoonish. While there are terms that labeled that
will eventually be canceled out, we need them as place holders so they are
initially defined.

{\bf Describe}\\
Knowns:\\
elevation, [m]: $\Delta \varepsilon$ \\
upper crustal density, [kg~m$^{-3}$]: $\rho_{t}$ \\
density contrast mantle to crust, [kg~m$^{-3}$]: $\Delta \rho = \rho_{m} -
\rho_{c}$ \\
gravitational acceleration, [m s$^{-2}$]: $g$

Unknowns:\\
height of column 1, [m]: $h_{1}$ \\
height of column 2, [m]: $h_{2}$ \\
thickness of mantle root, [m]: $r$ \\

Solve for $r$, [m].

{\bf Derive}\\
We start by equating columns pressures at a common depth of compensation,
\begin{equation}
    \underbrace{\rho_{c} g h_{1} + \rho_{m} g r}_{\mbox{\scriptsize column 1}} = 
    \underbrace{\rho_{t} g \Delta \varepsilon + \rho_{c} g (h_{2} - \Delta
    \varepsilon)}_{\mbox{\scriptsize column 2}}.
    \label{eq-column}
\end{equation}
Because $g$ is in every term in Equation~\ref{eq-column}, it can be canceled
out,
\begin{equation}
    \rho_{c} h_{1} + \rho_{m} r = 
    \rho_{t} \Delta \varepsilon + \rho_{c} (h_{2} - \Delta \varepsilon).
    \label{eq-column2}
\end{equation}

\begin{framed}
{\it Assumption:} The density of air is negligible so we do not need a pressure
term to account for the additional pressure of the air column.  To prove that we
don't need an air term, we can do a back of the envelope calculation to compare
the relative pressures of a lithospheric column to that of the air column.  The
density of air is on order 2~kg~m$^{-3}$ and the density of rock is on the order
of 3000~kg~m$^{-3}$.  The height of the air column is on the order of 1~km and
the height of the crust on the order of 30~km.  Hence, the ratio of the two
terms is
\begin{align*}
    \frac{\rho_c h_c}{\rho_{\text{air}} \Delta \varepsilon} &=
    \frac{(3000\text{ kg m$^{-3}$})(30\text{ km})}
    {(2\text{ kg m$^{-3}$})(1\text{ km})} \\
    &= \frac{90000}{2}\\
    &= 4.5 \times 10^{4}. \\
\end{align*}
It really doesn't matter if the heights vary considerably, the end result is
similar, the pressure due to a column of crust is on the order of 10,000 times
larger, so we can effectively ignore the influence of air above the column.
\end{framed}

Next, we equate column heights
\begin{equation}
    h_{2} = \Delta \varepsilon + h_{1} + r .
    \label{eq-height}
\end{equation}
Next, we can substitute Equation \ref{eq-height} into the modified pressure
Equation \ref{eq-column2}.  And after some algebra, we arrive at a solution for
the root,
\begin{align*}
    \rho_{c} h_{1} + \rho_{m} r &= \rho_{t} \Delta \varepsilon
        + \rho_{c}(\Delta \varepsilon + h_{1} + r - \Delta \varepsilon), \\
    \rho_{m} r - \rho_{c} r &= \rho_{t} \Delta \varepsilon 
        + \rho_{c} h_{1} - \rho_{c} h_{1}, \\
    (\rho_{m} - \rho_{c}) r &= \rho_{t} \Delta \varepsilon, \\
    r &= \frac{\rho_{t}}{\rho_{m} - \rho_{c}} \Delta \varepsilon. \\
\end{align*}
Recall from the problem statement that we are looking for a solution in terms of
$\Delta \rho$,
\begin{equation*}
\boxed{
    r = \frac{\rho_{t}}{\Delta \rho} \Delta \varepsilon
}\; .
\end{equation*}

{\bf Check}\\
Let us check the units,
\begin{align*}
    [\text{m}] &\qeq \frac{[\text{kg m}^{-3}]}{[\text{kg m}^{-3}]} [\text{m}], \\
    [\text{m}] &= [\text{m}].
\end{align*}

Our units check out.  Unfortunately, for this derivation there aren't any
additional tests we can use to verify the validity of the solution.

\end{document}
